\documentclass[../main]{subfiles}
\begin{document}
\ifSubfilesClassLoaded{\setcounter{page}{47}}{}
\section{Коммунизм и личность}
\label{sec:08_kommunizm_lichnost}

Проблема личности – одна из центральных проблем коммунизма. Значение личности при коммунизме возрастает до невиданных до сих пор высот. «На место старого   буржуазного    общества    с    его   классами     и  классовыми противоположностями приходит ассоциация, в которой свободное развитие каждого является условием свободного развития всех» - писали Маркс и Энгельс в «Манифесте коммунистической партии».

Всю предшествующую историю нужно понимать не просто как определенную последовательность событий, а как процесс становления сущности человека. Тогда становится ясно, что это в тоже время процесс становления определенного типа личности, отвечающей сущности человека. Какой должна быть эта личность? Этот вопрос каждый раз будет решаться практически и конкретно- исторически, но в любом случае это личность, сознательно творящая историю в коллективе таких же личностей-творцов. В свободной ассоциации людей вся их жизнь зависит от творческой деятельности, которая является наивысшим проявлением личности, её осуществляющей.

Личность – явление социальной природы от начала и до конца. Личностью индивид не рождается, а становится, включаясь в общественные отношения как активно действующий субъект. Первоначально человеческая жизнь задается новорожденному извне – из социального пространства, в котором он живет, со всеми вариациями отношений человека к человеку и человека к природе. И только освоив, то есть сделав своими, формы деятельности или, что то же самое, формы культуры, которые никоим образом в индивиде изначально не заложены, он становится личностью. Эти формы каждый, кто появляется на свет, застает уже готовыми. Для биологического организма они являются внешними, но только для организма, а не для личности. Мир деятельности как мир человеческой    культуры     –    это  действительный   мир     личности.   В коммунистическом обществе этот мир каждый должен сделать своим, а это значит - научиться изменять его.

Коммунистическая личность в своей практике будет активно творить общественные отношения, а не просто быть их пассивным продуктом. Проблема состоит в том, чтобы создать такие условия, в которых в своем индивидуальном развитии личность прошла бы в кратчайшие сроки путь становления общества и вышла на передний край общественной жизни. Нужно каждый раз искать и находить конкретные формы коллективности, которые бы связывали личность с обществом таким образом, чтобы она не могла бы не быть его активным членом и не жить всей полнотой общественной жизни. Поскольку общество постоянно развивается, такие формы «совместно-раздельной деятельности» не могут быть найдены один раз и навсегда. Управление общественным развитием при коммунизме состоит, главным образом, в том, чтобы научиться вовремя менять формы связи человека с обществом таким образом, чтобы они заставляли его действовать по-человечески в любых новых конкретно-исторических условиях.

Коммунистам вменяют в вину, то, что они хотят стандартизировать личность, сделать людей «одинаковыми». В этом обвинении сторонники капитализма, на самом деле выдают тайну современного капиталистического общества. Индивиды при капитализме в своей массе получаются настолько стандартизированными, что их индивидуальность определяется какими-то пустяками. Эта дурная индивидуальность не имеет отношения к истинной индивидуальности, как индивидуальному выражению освоенной человеком всеобщей человеческой культуры. Капитализм включает людей в стандартные практики производства и потребления, стандартизирует их проведение как во время работы, так и вне её. Люди уже «одинаковые». Коммунизм же предполагает   создание   общественного    пространства   для    развития индивидуальности.

Коммунисты видят необходимость в создании общества универсально- развитых личностей-творцов. В таком обществе индивидуальность каждого будет деятельным выражением и воплощением общественного развития. В каком смысле здесь можно говорить о стандартах?

Очевидно, человек не может быть универсальным творцом, не освоив необходимого для этого минимума человеческой культуры, и, соответственно, следует говорить о минимуме условий, который общество должно обеспечить каждому своему члену, чтобы он мог быть универсальным творцом, должен быть минимальный стандарт этих условий. В чем состоит этот самый минимум?

Сюда входят: определенный уровень развития системы здравоохранения, и именно направленной на сохранения здоровья и создания условий для здоровой жизни каждого человека, а не просто лечения болезней; качество и уровень образования, обязательный для каждого; система общественного питания и вообще всех общедоступных служб, которые могут уничтожить быт; обеспечение детства и старости. В одной из своих работ Ленин писал, что социализм начинается со стакана молока для каждого ребенка. Ясно, что наличие этого стакана еще не говорит о том, что общество развивается в направлении коммунизма, но вот без того, чтобы дети были сытыми, оно точно развиваться в этом направлении не может. Так же, как и не может развиваться без всего вышеперечисленного (список, конечно, не полный. Его разработкой должны заниматься люди, разбирающиеся в разных областях и по мере развития общества он будет расширяться).

Этот минимум человечного отношения к человеку не может быть раз и навсегда установившимся. Он изменяется в зависимости от условий, в которых происходит переход: от уровня развития человеческой культуры. В современный минимум, который был бы толчком для развития общества к коммунизму, обязательно должно входить высшее образование, включающее не только компьютерную грамотность на уровне программирования, но и научные основы понимания всех производственных процессов - политехнизм (для всех!). Также обязательным является обеспечение доступа ко всем современным средствам коммуникации.

Важно отметить, что даже самые развитые капиталистические страны не могут похвастаться поголовной грамотностью всего населения и доступностью медицинского обслуживания.

Необходимый минимум человечности в отношении к человеку, как уже говорилось, вещь историческая. Но этот минимум также должен включать необходимость его постоянного расширения. Чтобы коммунистическое общество, где свободное развитие каждого является условием развития всех, могло существовать и развиваться, в этот минимум должны постоянно включаться все новые передовые достижения человечества. Коммунистическое общество будет состоять из умных, здоровых и талантливых людей, которые, в конце концов, только и могут организовать общественное воспроизводство своей жизни на разумных началах.
\end{document}
